% sections/02_related_literature.tex
% Thematic (not paper-by-paper). Compressed ~1/3 vs the old three-subsection version.
% Empirical adoption-elasticity block moved up from Section 4 (referee m2).
% "Setup./Results." scaffolding and the coauthor asides are removed.
%
% Expected labels: app:taste sec:calibration
% Bibitem keys: ARS2024 Bayer2026 Langot2026 DCEGM2017 GallagherMuehlegger2011 Chandra2010
%               MuehleggerRapson2022 BeresteanuLi2011 Rust1987

\section{Related literature}\label{sec:lit}

We draw on two bodies of work: the macroeconomics of energy-shock stabilisation in
heterogeneous-agent open economies, and the microeconometrics of durable adoption
that disciplines the margin we add.

\subsection{Energy-shock stabilisation in heterogeneous-agent economies}

The closest work studies fiscal and monetary responses to the 2021--23 energy shock
in open HANK economies. \citet{ARS2024} provide the aggregate frame we inherit---an
energy-importing HANK with a CES energy nest, sticky import prices, and
intertemporal energy suppliers---and show that a low substitution elasticity makes
the shock contractionary through the real-income channel, so a subsidy stabilises
output, inflation and inequality at once; its drawback in their world is a
cross-country externality, since coordinated subsidies spike the world price.
\citet{Bayer2026} build a two-country monetary union to study exactly that
spillover, with a richer medium-scale block---two assets, capital, energy in
production---and a distributional mechanism that runs through within-country
energy-share heterogeneity: poorer households spend more on energy and suffer larger
real-income losses. \citet{Langot2026} take France as a small economy inside the
euro area, fed by the government's official energy-price forecasts, with a higher
intensive substitution elasticity and a subsistence energy floor that makes demand
non-homothetic; on a conditional-forecast path they find the French shield
effective, largely through a labour-cost channel that contains the price--wage
spiral.

These models disagree on the sign of the shield's welfare effect, and the
disagreement has an axis. Under perfectly elastic supply \citep{Langot2026} the
shield is effective; under a fixed union-wide energy quantity \citep{Bayer2026}
price suppression is destructive, because it blocks the substitution scarcity
requires. Our two supply closures---elastic price-taking and fixed-quantity---
reproduce the sign each obtains, which locates one source of the contrast in a
single parameter, the supply elasticity at the relevant horizon. We do not
over-read this. With no second country we cannot represent Bayer's cross-border
transmission, and our wage block, lacking an investment or employment margin,
reproduces the sign of Langot's pro-shield result but not its full labour-cost
mechanism.

What none of these models contains is an endogenous energy technology. ARS, Bayer,
and Langot hold the household's energy type fixed, so the shield's cost to the
transition is invisible to them. Our contribution is to make that type a choice.
The adoption margin adds a cost of price suppression the aggregate literature
cannot see---capping the brown price switches the transition off---without
displacing the ARS real-income channel, which operates in the background of every
experiment. Our heterogeneity is over discount factors and the adoption margin
rather than over energy shares, so we are silent on the energy-share-of-the-poor
mechanism; the two distributional channels are complements, not substitutes.

\subsection{Durable adoption and its identification}

The adoption margin is a discrete durable choice inside an incomplete-markets
household. We solve it with the discrete--continuous endogenous grid method of
\citet{DCEGM2017}, whose extreme-value taste shocks smooth the choice into
differentiable Jacobians for the sequence-space solution. In our calibration those
taste shocks are not only a numerical device: at a rare adoption margin the logit
scale is the structural semi-elasticity of switching
(Appendix~\ref{app:taste}), so it must be identified, not merely set small.

We discipline it against the quasi-experimental literature on durable adoption.
\citet{GallagherMuehlegger2011} and \citet{Chandra2010} estimate that a \$1{,}000 tax
incentive raises hybrid-vehicle adoption by 31--38\% among higher-income early
adopters; \citet{MuehleggerRapson2022} estimate a demand elasticity of $-2.1$ for
electric vehicles among low- and middle-income households, roughly 8\% per
\$1{,}000 at their mean transaction price, the mass-market floor on responsiveness.
\citet{BeresteanuLi2011} document the operating-cost margin---gasoline prices and
hybrid demand---the analogue of our energy-cost channel. These are estimates of the
upfront- and operating-cost elasticities the model needs, and
Section~\ref{sec:cal} maps them onto the taste scale. The structural
discrete-choice tradition that estimates such scales from rare, lumpy adjustment
\citep{Rust1987} is the right reference point for our regime, in contrast to the
interior-probability margins---labour-force participation---where the smoothing
reading applies and the scale is a numerical convenience.
